

\begin{abstract}
Indirect prompt-injection (IPI) benchmarks predominantly test planted
instructions that redirect a tool-using agent toward an attacker-chosen task.
Yet attacker-controlled content can instead tamper with the outcome of the
task already in progress. For example, a buyer's dispute evidence can cite a
false refund policy to inflate the refund without redirecting the agent away
from resolving the dispute. We define IPI by authority rather than
instruction-like wording: planted content constitutes an attack when it
attempts to influence a decision that its source is not authorized to
determine. We introduce a benchmark that makes this boundary explicit: in each
of 90 validated cases across three tool-use environments, the
user request specifies the source that resolves the targeted decision. Every
attack tampers with a decision required by the benign task rather than
replacing that task. Each case has attacks with the same harmful outcome that
target up to six parts of the agent's decision making: its objective, decision
rules, available actions, expected consequences, reading of observations, or
beliefs about the current state. Without a defense, attacks on three of the
five other targets (decision rules, available actions, and state) have higher
mean success than objective attacks across five LLMs. A detection-and-removal defense reduces objective-targeting attacks to
2.4\% success but leaves observation- and
state-targeting attacks at 9.5\% and
12.1\%; this gap persists under detection-optimized
rewrites. A tool-policy defense still allows
15.4\% attack success, close to 17.0\%
without a defense, because these attacks generally retain expected tools and
manipulate values obtained during execution. Evaluations whose attacks
redirect the task or target only the agent's objective can therefore
overestimate the coverage of these defenses. Code and data are available \mbox{at {\small\url{https://github.com/amazon-science/task-tampering-prompt-injection}}}.
\end{abstract}
